\documentclass[11pt]{scrartcl}

\usepackage[stixtwo]{fontsetup}
\usepackage{amsmath}

\begin{document}
\begin{enumerate}
	\item A nevező gyöktényezős alakban:
	      \[
		      x^2 - 5x + 6 = (x - 3)(x - 2)
		      \text.
	      \]
	      A nullával való osztást nem értelmezzük, ezért $D_f = \mathbb{R}
		      \setminus \{2, 3\}$.

	      A függvény $D_f$ minden pontjában folytonos, mert folytonos függvények
	      hányadosa.

	\item Az $x = 2$ pontban:
	      \[
		      \lim_{x \to 2} f(x)
		      = \lim_{x \to 2}
		      \underbrace{\frac{\sin\left(x^2 - 4\right)}{x^2 - 4}}_{\to 1} \cdot
		      \underbrace{\frac{x (x + 2)}{x - 3}}_{\to \frac{2 \cdot 4}{-1}}
		      = -8,
	      \]
	      a bal és a jobb oldali határérték is $-8$, így itt megszüntethető
	      (elsőfajú) szakadása van a függvénynek.

	      Az $x = 3$ pontban:
	      \[
		      \lim_{x \to 3 \pm 0} f(x)
		      = \lim_{x \to 3 \pm 0}
		      \frac{
			      \overbrace{x \sin\left(x^2 - 4\right)}^{\to 3 \sin 5 < 0}
		      }{
			      \underbrace{(x - 2)(x - 3)}_{\to 0 \pm 0}
		      }
		      = \mp \infty,
	      \]
	      vagyis a bal oldali határérték $+\infty$, a jobb oldali $-\infty$, így
	      itt másodfajú szakadása van a függvénynek.
\end{enumerate}
\end{document}
